Integrating this measure along physiological paths yields a
parameter-free gene sensitivity metric, $\kmu$: each enzyme's
rerouting burden. In carbon-starved \emph{E.~coli}, $\kmu$
predicts transcriptional log-fold changes across $424$ genes
($r = +0.395$, $p = 2.6 \times 10^{-17}$; partial $r = +0.269$),
robust across tie-breaking rules and stress axes. The induced
genes sit in the operons of the global carbon and energy regulons,
and the rerouting mass concentrates on the fork metabolites of
central carbon. The association vanishes at the protein layer
($r = -0.083$, $366$ genes, matched quantitative proteomics):
cells transcribe standby capacity for rerouting, buffering protein
synthesis, so the persistent memory of cyclic perturbations
($66\%$ non-reverting) must live in fast post-translational state.
Double-knockout epistasis mirrors active-set boundary overlaps
($\rho_S = 0.865$). Our framework unites linear programming,
discrete geometry, and transcriptional regulation into an
accessible, predictive foundation for metabolic systems biology.
